\documentclass[12pt,a4paper]{article}
\usepackage[margin=1in]{geometry}
\usepackage[T1]{fontenc}
\usepackage{setspace}
\usepackage{amsmath,amssymb}
\usepackage{booktabs}
\usepackage{longtable}
\usepackage{xcolor}
\usepackage{hyperref}
\usepackage{fancyhdr}

\setstretch{1.18}
\setlength{\parskip}{0.45em}
\setlength{\parindent}{0pt}

\hypersetup{colorlinks=true,linkcolor=blue,citecolor=blue,urlcolor=blue}

\pagestyle{fancy}
\fancyhf{}
\lhead{CST428 - Blockchain Technologies}
\rhead{Module 1 Master Notes}
\cfoot{\thepage}

\newcommand{\examtip}[1]{%
\vspace{0.35em}
\noindent\fcolorbox{black}{gray!12}{\parbox{0.965\linewidth}{\textbf{Exam Tip:} #1}}%
\vspace{0.35em}
}

\title{\textbf{CST428 Blockchain Technologies}\\\Large Module 1 High-Clarity Notes\\\normalsize Theory + Hard Topic Explanations + Exam-Ready Answers}
\author{Prepared from Module 1 source PDFs and all available previous-year Module 1 questions}
\date{\today}

\begin{document}

\section{Module 1 Topic Map From Previous Questions}
From all available question papers (May 2024, June 2023, October 2023), Module 1 repeatedly asks:
\begin{itemize}
\item Symmetric vs asymmetric cryptography
\item AES fundamentals and state matrix flow
\item RSA algorithm and RSA numericals
\item Elliptic Curve Cryptography (ECC)
\item Hash function properties
\item SHA-256 algorithm and compression function
\item Merkle trees (benefits and blockchain use)
\item Distributed Hash Tables (DHT)
\end{itemize}

\section{Core Theory (Explained Clearly)}

\subsection{Introduction to Cryptography}
Cryptography is the science of protecting information when adversaries are present.

Basic terms:
\begin{itemize}
\item \textbf{Plaintext (P)}: original readable message.
\item \textbf{Ciphertext (C)}: encrypted unreadable output.
\item \textbf{Encryption (E)}: process converting plaintext to ciphertext.
\item \textbf{Decryption (D)}: process recovering plaintext from ciphertext.
\item \textbf{Key (K)}: secret/control parameter used by algorithms.
\end{itemize}

Security goals:
\begin{itemize}
\item \textbf{Confidentiality:} unauthorized users must not read data.
\item \textbf{Integrity:} unauthorized changes must be detectable.
\item \textbf{Authentication:} identity or source must be verifiable.
\item \textbf{Non-repudiation:} sender cannot deny performed action.
\item \textbf{Accountability:} actions can be traced to responsible entities.
\end{itemize}

\subsection{Symmetric vs Asymmetric Cryptography}
\begin{longtable}{p{0.31\textwidth}p{0.62\textwidth}}
\toprule
\textbf{Symmetric Cryptography} & \textbf{Asymmetric Cryptography} \\
\midrule
One key is used for both encryption and decryption & Two keys are used: public key and private key \\
Very fast and efficient for large data & Slower, used for secure key exchange and signatures \\
Key sharing is a challenge & Public key can be shared openly \\
Mainly confidentiality & Confidentiality + authentication + non-repudiation \\
Examples: AES, DES & Examples: RSA, ECC, DSA \\
\bottomrule
\end{longtable}

Mathematical form:
\begin{align}
\text{Symmetric: } & C = E_K(P), \quad P = D_K(C) \\
\text{Asymmetric: } & C = E_{K_{pub}}(P), \quad P = D_{K_{priv}}(C)
\end{align}

\subsection{AES (Advanced Encryption Standard)}
AES is a symmetric block cipher.

Key points:
\begin{itemize}
\item Block size: 128 bits.
\item Key sizes: 128, 192, 256 bits.
\item Number of rounds: 10, 12, 14 respectively.
\item Works on a \(4 \times 4\) state matrix of bytes.
\end{itemize}

Round operations:
\begin{enumerate}
\item \textbf{AddRoundKey}: XOR state with round key.
\item \textbf{SubBytes}: nonlinear substitution using S-box.
\item \textbf{ShiftRows}: row-wise cyclic shifts.
\item \textbf{MixColumns}: column-wise linear mixing.
\end{enumerate}

Final round has no MixColumns.

\examtip{In exam answers, always write the four AES transformations in order and explicitly mention that MixColumns is omitted in the final round.}

\subsection{RSA Cryptography}
RSA is based on the practical difficulty of factoring a large integer \(n = pq\), where \(p\) and \(q\) are prime.

\subsubsection{RSA Key Generation}
\begin{itemize}
\item Choose two large primes \(p\) and \(q\).
\item Compute modulus: \(n = pq\).
\item Compute Euler totient: \(\phi(n) = (p-1)(q-1)\).
\item Choose \(e\) such that \(1 < e < \phi(n)\) and \(\gcd(e,\phi(n))=1\).
\item Compute \(d\) such that:
\begin{equation}
ed \equiv 1 \pmod{\phi(n)}
\end{equation}
\item Public key is \((n,e)\), private key is \((n,d)\).
\end{itemize}

RSA operations:
\begin{align}
C &= M^e \bmod n \\
M &= C^d \bmod n
\end{align}

\subsection{Elliptic Curve Cryptography (ECC)}
ECC works over points on elliptic curves over finite fields.

Common curve form:
\begin{equation}
y^2 = x^3 + ax + b
\end{equation}

ECC idea:
\begin{itemize}
\item choose a base point \(G\) on curve,
\item private key is a random integer \(d\),
\item public key is \(Q=dG\).
\end{itemize}

Why ECC is important:
\begin{itemize}
\item same security with smaller key sizes compared to RSA,
\item lower storage and bandwidth,
\item better performance for constrained systems.
\end{itemize}

\subsection{Digital Signatures (RSA/ ECC Context)}
Digital signatures provide:
\begin{itemize}
\item data origin authentication,
\item integrity,
\item non-repudiation.
\end{itemize}

Typical signing workflow:
\begin{enumerate}
\item compute hash \(h(m)\),
\item sign hash with private key,
\item verify using public key.
\end{enumerate}

\subsection{Hash Functions and Properties}
A hash function maps arbitrary-length input to fixed-length digest.

Practical properties:
\begin{itemize}
\item fixed-length output,
\item fast computation.
\end{itemize}

Security properties:
\begin{itemize}
\item \textbf{Pre-image resistance}: from digest, infeasible to recover input.
\item \textbf{Second pre-image resistance}: for given input \(x\), infeasible to find \(x'\neq x\) with same hash.
\item \textbf{Collision resistance}: infeasible to find any two different inputs with same hash.
\end{itemize}

Avalanche effect: tiny input change causes major output change.

\subsection{SHA-256}
SHA-256 belongs to SHA-2.

Specifications:
\begin{itemize}
\item input block size: 512 bits,
\item word size: 32 bits,
\item digest size: 256 bits,
\item 64 rounds per block.
\end{itemize}

High-level steps:
\begin{itemize}
\item Padding and length appending.
\item Split input into 512-bit blocks.
\item Initialize 8 hash words from SHA-256 standard constants.
\item Prepare message schedule (80 words derived from block).
\item Run 64-round compression function for each block.
\item Accumulate state and produce 256-bit final digest.
\end{itemize}

Compression equations for each round \(t\):
\begin{align}
T_1 &= h + \Sigma_1(e) + Ch(e,f,g) + K_t + W_t \\
T_2 &= \Sigma_0(a) + Maj(a,b,c)
\end{align}
where
\begin{align}
Ch(x,y,z) &= (x \wedge y) \oplus (\neg x \wedge z) \\
Maj(x,y,z) &= (x \wedge y) \oplus (x \wedge z) \oplus (y \wedge z)
\end{align}

\subsection{Merkle Trees}
Merkle tree is a binary hash tree.

Structure:
\begin{itemize}
\item leaves: hash of transactions/data,
\item parent: hash(left child || right child),
\item top node: Merkle root.
\end{itemize}

Benefits:
\begin{itemize}
\item fast integrity verification,
\item efficient inclusion proofs,
\item tamper detection,
\item storage and bandwidth efficiency.
\end{itemize}

\subsection{Distributed Hash Tables (DHT)}
DHT is a distributed key-value lookup system.

Workflow:
\begin{enumerate}
\item hash the key to an identifier,
\item route request to node responsible for that identifier,
\item node stores or returns value,
\item replication improves availability and fault tolerance.
\end{enumerate}

Applications include P2P networks and decentralized data systems.

\section{Hard Topics Decoded (Numericals + Concepts)}

\subsection{How To Solve RSA Numericals Correctly}
Use this exact method in exams.

Given \(p,q,e,M\):
\begin{itemize}
\item Step 1: Compute \(n=pq\) - multiply the two primes to get modulus.
\item Step 2: Compute \(\phi(n)=(p-1)(q-1)\) - subtract 1 from each prime and multiply.
\item Step 3: Find \(d\) such that \(ed \equiv 1 \pmod{\phi(n)}\) - solve using extended Euclidean algorithm or trial.
\item Step 4: Encrypt using \(C=M^e \bmod n\) - raise message to public exponent, take modulo.
\item Step 5: Decrypt using \(M=C^d \bmod n\) - raise ciphertext to private exponent, take modulo.
\item Step 6: Verify recovered message equals original message to confirm correctness.
\end{itemize}

\textbf{Fast modular exponentiation tip:}
Break power into squares and reduce modulo at each step to avoid huge numbers.

\subsection{How To Explain AES State Matrix For Full Marks}
\begin{itemize}
\item AES takes 128-bit plaintext and arranges it into 16-byte state matrix.
\item State matrix is transformed round by round through the algorithm.
\item Start with AddRoundKey pre-round transformation using initial key.
\item For each normal round: SubBytes, ShiftRows, MixColumns, AddRoundKey in order.
\item Final round omits MixColumns step to save computation.
\item Final state matrix is read out as 16-byte ciphertext.
\end{itemize}

\subsection{How To Explain SHA-256 Compression Clearly}
\begin{itemize}
\item Mention eight working variables \((a,b,c,d,e,f,g,h)\).
\item Define round inputs \(W_t\) (message schedule) and \(K_t\) (round constants).
\item Write \(T_1\) and \(T_2\) formulas explicitly.
\item Show register update flow: \(h \leftarrow g\), \(g \leftarrow f\), \(f \leftarrow e\), etc.
\item Explain 64 rounds per block with different \(W_t\) and \(K_t\).
\item Explain final state addition and 256-bit output accumulation.
\end{itemize}

\subsection{How To Explain ECC Without Confusion}
\begin{itemize}
\item Start from curve equation and finite-field setting.
\item Define base point \(G\), private key \(d\), public key \(Q=dG\).
\item Explain hard problem: reversing scalar multiplication is infeasible.
\item Explain smaller key-size advantage over RSA for same security level.
\item Mention ECC is efficient for blockchain and constrained systems.
\item Conclude with security-efficiency tradeoff benefit.
\end{itemize}

\newpage
\section*{Solved Question Papers (Module 1)}
\addcontentsline{toc}{section}{Solved Question Papers (Module 1)}

\subsection*{Part A (3 Marks) - Every Answer Has 6 Points}

\subsubsection*{Q1 May 2024: Differentiate between symmetric and asymmetric cryptography}
\begin{itemize}
\item Symmetric uses one key for both encryption and decryption.
\item Asymmetric uses two keys: public key and private key.
\item Symmetric is faster and suitable for large data encryption.
\item Asymmetric is slower and mainly used for key exchange and signatures.
\item Symmetric has key distribution challenge between communicating users.
\item Asymmetric supports authentication and non-repudiation through key pairs.
\end{itemize}

\subsubsection*{Q2 May 2024: Benefits of Merkle trees}
\begin{itemize}
\item Merkle tree stores transaction hashes hierarchically in binary tree form.
\item Merkle root summarizes entire dataset integrity in single hash.
\item Any leaf change propagates upward and changes root hash immediately.
\item Inclusion proofs are efficient using only logarithmic proof path length.
\item Reduces storage because hashes are much smaller than full data blocks.
\item Useful for lightweight verification in blockchain clients without full blocks.
\end{itemize}

\subsubsection*{Q1 June 2023: Compare symmetric and asymmetric key cryptography}
\begin{itemize}
\item Symmetric key is shared secret between users; asymmetric has public-private key pair.
\item Symmetric encryption and decryption are computationally very fast.
\item Asymmetric operations are computationally heavier and slower.
\item Symmetric is ideal for bulk confidentiality of large data volumes.
\item Asymmetric is ideal for key distribution and digital signatures.
\item Asymmetric improves trust in open networks without pre-shared secrets.
\end{itemize}

\subsubsection*{Q2 June 2023: How secure hash functions strengthen blockchain}
\begin{itemize}
\item Hash links each block to previous block through cryptographic hash chain.
\item Hashes provide tamper evidence through integrity checks that fail on modification.
\item SHA-256 based computational puzzles enable Proof-of-Work mining security.
\item Merkle roots commit all transactions compactly in single 256-bit hash.
\item Collision resistance prevents attacker from easily substituting fake transaction data.
\item Avalanche effect makes manipulation immediately visible as digest changes drastically.
\end{itemize}

\subsubsection*{Q1 October 2023: Major differences between symmetric and asymmetric cryptography}
\begin{itemize}
\item Number of keys: one secret key (symmetric) vs two keys public-private pair (asymmetric).
\item Speed: symmetric is very fast, asymmetric is significantly slower.
\item Typical use: symmetric for bulk data protection, asymmetric for key setup and signatures.
\item Key sharing: symmetric needs secure channel; asymmetric can share public key openly.
\item Security services: asymmetric naturally provides authentication and non-repudiation.
\item Common examples: AES (symmetric), RSA and ECC (asymmetric).
\end{itemize}

\subsubsection*{Q2 October 2023: Describe distributed hash tables and their working}
\begin{itemize}
\item DHT is a distributed key-value lookup mechanism across many peers.
\item Key is hashed to produce an identifier in a logical keyspace.
\item Each node is responsible for specific identifier ranges (ownership shares).
\item Queries are routed through overlay neighbors to node owning that identifier.
\item Responsible node stores and retrieves the value for that key.
\item DHT provides decentralized, scalable, and fault-tolerant lookup without central server.
\end{itemize}

\subsection*{Part B (14 Marks) - Every Answer Has At Least 18 Points}

\subsubsection*{May 2024 Q11 (a+b): RSA concepts + hash properties + DHT working}
\textbf{Exam-ready 18 points:}
\begin{itemize}
\item RSA is an asymmetric algorithm based on integer factorization hardness.
\item Choose two large primes \(p\) and \(q\).
\item Compute modulus \(n=pq\).
\item Compute Euler totient \(\phi(n)=(p-1)(q-1)\).
\item Select public exponent \(e\) such that \(1<e<\phi(n)\) and co-prime condition holds.
\item Compute private exponent \(d\) from \(ed \equiv 1 \pmod{\phi(n)}\).
\item Public key is \((n,e)\) and private key is \((n,d)\).
\item RSA encryption formula is \(C=M^e \bmod n\).
\item RSA decryption formula is \(M=C^d \bmod n\).
\item RSA gives confidentiality; with signing it also gives authentication.
\item Hash function property 1: fixed-length digest for arbitrary input length.
\item Hash function property 2: fast computation.
\item Hash function property 3: pre-image resistance.
\item Hash function property 4: second pre-image resistance.
\item Hash function property 5: collision resistance.
\item DHT stores key-value pairs across multiple nodes, not one server.
\item DHT request is routed to node responsible for hashed key identifier.
\item DHT improves scalability and removes centralized lookup bottleneck.
\end{itemize}

\subsubsection*{May 2024 Q12 (a+b): ECC algorithm + SHA-256 compression}
\textbf{Exam-ready 18 points:}
\begin{itemize}
\item ECC is based on elliptic curve discrete logarithm difficulty.
\item Standard curve form is \(y^2=x^3+ax+b\) over finite field.
\item Choose domain parameters including base point \(G\) and order \(n\).
\item Select private key \(d\), with \(1\le d \le n-1\).
\item Compute public key \(Q=dG\).
\item Key exchange derives shared secret by scalar multiplication.
\item ECC needs smaller keys than RSA for similar security strength.
\item Smaller keys reduce storage and communication overhead.
\item ECC is efficient for resource-constrained environments.
\item SHA-256 processes input in 512-bit blocks.
\item SHA-256 has 8 working registers \((a,b,c,d,e,f,g,h)\).
\item Message schedule \(W_t\) and constant \(K_t\) are used each round.
\item Compute \(T_1 = h + \Sigma_1(e) + Ch(e,f,g) + K_t + W_t\).
\item Compute \(T_2 = \Sigma_0(a) + Maj(a,b,c)\).
\item Update registers and continue for 64 rounds.
\item Add final working values to previous hash state.
\item Repeat for all message blocks.
\item Concatenate final state to obtain 256-bit digest.
\end{itemize}

\subsubsection*{June 2023 Q11 (a+b): AES state matrix + RSA numerical}
\textbf{Exam-ready 18 points:}
\begin{itemize}
\item AES plaintext block size is 128 bits (16 bytes).
\item Arrange 16 bytes into a \(4\times4\) state matrix.
\item AES key sizes are 128, 192, and 256 bits.
\item Round counts are 10, 12, and 14 respectively.
\item Initial step is AddRoundKey with round-0 key.
\item Normal rounds apply SubBytes transformation (S-box substitution).
\item Then apply ShiftRows transformation (cyclic row shifts).
\item Then apply MixColumns transformation (linear column mixing).
\item Then apply AddRoundKey again (XOR with round key).
\item Final round omits MixColumns step.

\textbf{RSA Numerical: } \(p=7, q=11, e=17, M=25\)
\item Compute modulus: \(n = p \times q = 7 \times 11 = 77\)
\item Compute totient: \(\phi(n) = (p-1)(q-1) = (7-1)(11-1) = 6 \times 10 = 60\)
\item Given \(e=17\), find \(d\) from \(ed \equiv 1 \pmod{60}\)
\item Test: \(17 \times 3 = 51 \equiv 51 \pmod{60}\) (no), \(17 \times 53 = 901 = 15 \times 60 + 1 \equiv 1 \pmod{60}\) (yes!), so \(d=53\)
\item Public key: \((n,e)=(77,17)\), Private key: \((n,d)=(77,53)\)
\item Encrypt \(M=25\): \(C = 25^{17} \bmod 77\). Using successive squaring: \(25^2=625 \equiv 16 \pmod{77}\), \(25^4 \equiv 256 \equiv 25 \pmod{77}\), \(25^{16} \equiv 25 \pmod{77}\), so \(C = 25 \times 16 \times 25 \equiv 9 \pmod{77}\)
\item Decrypt \(C=9\): \(M = 9^{53} \bmod 77 = 25\)
\item Recovered message equals original, confirming RSA correctness.
\end{itemize}

\subsubsection*{June 2023 Q12 (a+b): SHA-256 + Merkle tree security role}
\textbf{Exam-ready 18 points:}
\begin{itemize}
\item SHA-256 belongs to SHA-2 family.
\item Input is padded to satisfy block alignment rules.
\item Message is split into 512-bit blocks.
\item Initial hash state consists of 8 fixed 32-bit words.
\item For each block, create message schedule words.
\item Perform 64 compression rounds.
\item Use functions \(Ch\), \(Maj\), \(\Sigma_0\), and \(\Sigma_1\).
\item Use constants \(K_t\) and schedule \(W_t\).
\item Update eight working registers each round.
\item Add working registers to hash state after rounds.
\item Repeat across all blocks.
\item Final output length is 256 bits.
\item Merkle tree uses hashes at leaves and internal nodes.
\item Parent node is hash of concatenated child hashes.
\item Root hash commits whole transaction set.
\item Inclusion proof requires only sibling hashes along one path.
\item Any transaction change alters root hash and reveals tampering.
\item Therefore Merkle tree strengthens integrity and lightweight verification.
\end{itemize}

\subsubsection*{October 2023 Q11 (a+b): hash properties + public/private keys + RSA numerical}
\textbf{Exam-ready 18 points:}
\begin{itemize}
\item Hash functions convert variable-length input to fixed-length digest.
\item Hash functions are one-way in practical cryptography.
\item Security property 1: pre-image resistance (cannot reverse hash).
\item Security property 2: second pre-image resistance (cannot find alternate input).
\item Security property 3: collision resistance (cannot find two inputs with same hash).
\item Additional property: fast deterministic computation.
\item In blockchain, hashes link blocks and cryptographically preserve integrity.
\item Public key is openly distributable to any party.
\item Private key must be kept secret by the owner.
\item In encryption, receiver public key is used to encipher.
\item In decryption, receiver private key is used to decipher.

\textbf{RSA Numerical: } \(p=3, q=11, e=7, M=5\)
\item Compute modulus: \(n = p \times q = 3 \times 11 = 33\)
\item Compute totient: \(\phi(n) = (p-1)(q-1) = (3-1)(11-1) = 2 \times 10 = 20\)
\item Given \(e=7\), find \(d\) from \(ed \equiv 1 \pmod{20}\)
\item Test: \(7 \times 3 = 21 = 1 \times 20 + 1 \equiv 1 \pmod{20}\), so \(d=3\)
\item Public key: \((n,e)=(33,7)\), Private key: \((n,d)=(33,3)\)
\item Encrypt \(M=5\): \(C = 5^7 \bmod 33 = 78125 \bmod 33\). Note \(78125 = 2367 \times 33 + 14\), so \(C=14\)
\item Decrypt \(C=14\): \(M = 14^3 \bmod 33 = 2744\). Note \(2744 = 83 \times 33 + 5\), so \(M=5\)
\item Recovery of original message validates the RSA process.
\end{itemize}

\subsubsection*{October 2023 Q12 (a+b): Merkle tree with neat explanation + SHA-256 algorithm}
\textbf{Exam-ready 18 points:}
\begin{itemize}
\item Merkle tree is a binary tree of hashes.
\item Leaf nodes store hashes of transactions or data blocks.
\item Internal node value is hash(left child \(\|\) right child) where \(\|\) is concatenation.
\item Top node is the Merkle root stored in block header.
\item Root commitment allows verification of any transaction inclusion.
\item Inclusion proofs use only branch hashes, not full transaction data.
\item This reduces bandwidth and verification cost significantly.
\item Any data tampering at leaf level changes leaf hash immediately.
\item Changed leaf hash propagates upward and changes root hash.
\item Hence Merkle trees provide efficient integrity checking.
\item SHA-256 starts with preprocessing and padding step.
\item Message is divided into 512-bit blocks.
\item Eight 32-bit hash variables are initialized to fixed values.
\item Message schedule expands data for each round use.
\item Each block runs 64 compression rounds sequentially.
\item Round formulas compute \(T_1\) and \(T_2\) from current state.
\item Registers are updated in each round: \(h\leftarrow g, g\leftarrow f, \ldots, a\leftarrow T_1+T_2\).
\item Final state accumulation produces 256-bit digest.
\end{itemize}

\section{Diagrams and Visual References}

\subsection{RSA Key Generation and Use Diagram}
\vspace{0.3em}
\noindent\textbf{RSA Process Flow:}
\begin{center}
\framebox[0.9\linewidth][c]{
  \begin{minipage}{0.85\linewidth}
    \textbf{Step 1: Key Generation}\\
    Select primes $p$, $q$ $\rightarrow$ Compute $n = pq$, $\phi(n) = (p-1)(q-1)$\\
    Choose $e$ with $\gcd(e, \phi(n)) = 1$ $\rightarrow$ Find $d$ such that $ed \equiv 1 \pmod{\phi(n)}$\\
    \textbf{Public Key:} $(n, e)$ \quad \textbf{Private Key:} $(n, d)$\
    \\[0.2cm]
    \textbf{Step 2: Encryption (Sender)}\\[0.1cm]
    Message $M$ $\rightarrow$ $C = M^e \bmod n$ $\rightarrow$ Send ciphertext $C$\
    \\[0.2cm]
    \textbf{Step 3: Decryption (Receiver)}\\[0.1cm]
    Receive $C$ $\rightarrow$ $M = C^d \bmod n$ $\rightarrow$ Plaintext $M$
  \end{minipage}
}
\end{center}
\vspace{0.3em}

\subsection{Merkle Tree Structure Diagram}
\vspace{0.3em}
\noindent\textbf{4-Level Merkle Tree Hierarchy:}
\begin{center}
\framebox[0.95\linewidth][c]{
  \begin{minipage}{0.90\linewidth}
    \vspace{0.15cm}
    \centerline{\textbf{Root} = $H(H_{AB} || H_{CD})$}
    \vspace{0.1cm}
    \centerline{$\Big\downarrow$}
    \vspace{0.05cm}
    \centerline{$H_{AB} = H(H_A || H_B)$ \hspace{2cm} $H_{CD} = H(H_C || H_D)$}
    \vspace{0.1cm}
    \centerline{$\Big\downarrow$ \hspace{2cm} $\Big\downarrow$ \hspace{1.5cm} $\Big\downarrow$ \hspace{2cm} $\Big\downarrow$}
    \vspace{0.05cm}
    \centerline{$H_A$ \hspace{1.5cm} $H_B$ \hspace{1.5cm} $H_C$ \hspace{1.5cm} $H_D$}
    \vspace{0.1cm}
    \centerline{$\Big\downarrow$ \hspace{1.5cm} $\Big\downarrow$ \hspace{1.5cm} $\Big\downarrow$ \hspace{1.5cm} $\Big\downarrow$}
    \vspace{0.05cm}
    \centerline{Tx-A \hspace{1.5cm} Tx-B \hspace{1.5cm} Tx-C \hspace{1.5cm} Tx-D}
    \vspace{0.1cm}
    \textbf{Property:} Any change in one transaction changes root hash $\Rightarrow$ Immutability detection.
    \vspace{0.1cm}
  \end{minipage}
}
\end{center}
\vspace{0.3em}

\subsection{AES Encryption Flow Diagram}
\vspace{0.3em}
\noindent\textbf{AES Block Encryption Pipeline (128-bit key, 10 rounds):}
\begin{center}
\framebox[0.95\linewidth][c]{
  \begin{minipage}{0.90\linewidth}
    \vspace{0.1cm}
    \centerline{128-bit Plaintext}
    \vspace{0.08cm}
    \centerline{$\Big\downarrow$}
    \vspace{0.05cm}
    \centerline{Parse into 4×4 State Matrix (16 bytes)}
    \vspace{0.08cm}
    \centerline{$\Big\downarrow$}
    \vspace{0.05cm}
    \centerline{Initial AddRoundKey (XOR with Round Key 0)}
    \vspace{0.08cm}
    \centerline{$\Big\downarrow$}
    \vspace{0.1cm}
    \centerline{\textbf{Main Rounds (10 times for 128-bit key):}}
    \vspace{0.05cm}
    \centerline{SubBytes $\rightarrow$ ShiftRows $\rightarrow$ MixColumns $\rightarrow$ AddRoundKey}
    \vspace{0.08cm}
    \centerline{$\Big\downarrow$}
    \vspace{0.1cm}
    \centerline{\textbf{Final Round (omit MixColumns):}}
    \vspace{0.05cm}
    \centerline{SubBytes $\rightarrow$ ShiftRows $\rightarrow$ AddRoundKey}
    \vspace{0.08cm}
    \centerline{$\Big\downarrow$}
    \vspace{0.05cm}
    \centerline{128-bit Ciphertext}
    \vspace{0.1cm}
  \end{minipage}
}
\end{center}
\vspace{0.3em}

\section{Quick Revision Sheets}
\subsection{Formula Sheet}
\begin{align}
\phi(n) &= (p-1)(q-1) \\
ed &\equiv 1 \pmod{\phi(n)} \\
C &= M^e \bmod n \\
M &= C^d \bmod n \\
T_1 &= h + \Sigma_1(e) + Ch(e,f,g) + K_t + W_t \\
T_2 &= \Sigma_0(a) + Maj(a,b,c)
\end{align}

\subsection{One-Page Answer Pattern}
\begin{itemize}
\item Start with 1-line definition of the topic.
\item Write algorithm or structure in ordered steps or bullets.
\item Add equations with symbols and their meanings.
\item Add one practical blockchain relevance line.
\item End with one-line conclusion linked to security goal.
\end{itemize}

\examtip{If you run short on time: for 14-mark answers, first write all 18 point headers quickly, then fill equations and one-line explanations. Structured answers fetch marks even under time pressure.}

\section*{Closing Note}
\addcontentsline{toc}{section}{Closing Note}
This rewritten Module 1 version is optimized for understanding plus exam writing precision. Follow the hard-topic sections first, then practice reproducing the solved answers with the same point structure.

\end{document}
